\chapter{Limits and conclusion}\label{ch:end}

\section{Limits}
\paragraph{The test set is small} Fifty-two requests give intervals about seven points wide. The ranking of the deployed tagger over the HMM and CRF is clear; its lead over the rules alone is not yet significant.

\paragraph{The training data is synthetic} Templates and catalogue titles do not sound like customers. The gazetteer covers for this today. Real requests from the application, with consent, are the remedy, and the plan's benchmark is the measure.

\paragraph{Pictures are read coarsely} Zero-shot CLIP is reliable for garment type, colour and print on whole pictures, and not for necklines or sleeves. Attribute heads trained on fashion images, as the plan proposes with FashionSigLIP, are needed before the board can read fine detail.

\paragraph{Generated pictures can be wrong} Image models draw necklines and sleeves wrongly often enough that every published picture is reviewed by hand. The ranking narrows the choice; it does not replace the review.

\paragraph{Ease values are defaults} The allowance added for each fit is a starting value, marked as such on the sheet. Partner tailors' own values will replace it.

\section{Conclusion}
Bespoke shows that the gap between a customer's loose, mixed-language description and an instruction a tailor can cut from can be closed with a small, careful system that runs in the browser. The ontology carries the domain knowledge; a gazetteer and a fine-tuned 33-million-parameter encoder together tag requests with F1 0.953; a binder and a fold turn spans into a design whose every value has a source; and a question policy asks only what a tailor needs, in the order that matters. A language model is used where it helps and is never trusted unchecked. Choices are made by looking, from pictures made on a laptop with an open model. The application is live, tested and fast. What remains is scale: the data platform that learns from orders, and the tailors who will judge whether the sheet is enough.
